\subsection*{3.2\; The four species, written out}

\textbf{This is the level the model actually reaches, so it is worth writing without abbreviation.}
Everything above this point is mechanics: a surface with a strain energy, an adhesion with a rest
length, a tissue that grows. Proteolysis is the first place the model is a \emph{chemistry} --- five
molecular species, four of them proteins that the literature names, with concentrations that obey
transport equations and a rate law taken from an assay rather than fitted to a picture. The claim
being made by writing them out is a specific one: the breach in these runs is not a hole painted onto
a mesh, it is the consequence of \emph{molecules being made, moving, binding one another, and cutting
collagen~IV}, and every symbol below has a counterpart in a paper.

\smallskip
{\footnotesize
\begin{tabular}{@{}lllll@{}}\toprule
symbol & species & UniProt / EC & where it lives & transported?\\\midrule
$T$      & MT1-MMP (MMP-14)   & P50281 & per-\code{cell} state, mapped to \code{bm\_face} & \textbf{no} --- transmembrane\\
$c_P$    & proMMP-2, zymogen  & P08253 & \textbf{field} on \code{bm\_face} & yes, $D_M$\\
$c_M$    & MMP-2, active      & EC 3.4.24.24 & \textbf{field} on \code{bm\_face} & yes, $D_M$\\
$c_T$    & TIMP-2             & P16035 & \textbf{field} on \code{bm\_face} & yes, $D_T$\\
$c_3$    & TIMP-3             & P35625 & per-\code{bm\_face} state & \textbf{no} --- ECM-bound\\
\bottomrule\end{tabular}}
\smallskip

\textbf{The four equations.} All are written on the sheet's own surface, so $\nabla_s^2$ is the
Laplace--Beltrami operator of the current \emph{deformed} triangulation and not of a flat chart; $s_P$,
$s_T$, $s_3$ and $T$ are per-cell fields carried onto the faces each cell touches, so the epithelium
decides where its enzymes go. Writing $x=(c_T+c_3)/K$ for the inhibitor's occupancy of the tethered
enzyme:
%
\begin{align}
\frac{\partial c_P}{\partial t}
  &= \underbrace{D_M\nabla_s^2 c_P}_{\text{diffusion}}
   + \underbrace{s_P(\mathbf x)}_{\text{secretion}}
   - \underbrace{a(\mathbf x,t)}_{\text{activation}}
   - \underbrace{c_P/\tau_P}_{\text{clearance}} \label{eq:pro}\\[2pt]
\frac{\partial c_M}{\partial t}
  &= D_M\nabla_s^2 c_M
   + a(\mathbf x,t)
   - \underbrace{k_i\,c_M\,(c_T+c_3)}_{\text{1:1 inhibition}}
   - c_M/\tau_M \label{eq:mmp}\\[2pt]
\frac{\partial c_T}{\partial t}
  &= D_T\nabla_s^2 c_T + s_T(\mathbf x) - k_i\,c_M\,(c_T+c_3) - c_T/\tau_T \label{eq:timp}\\[2pt]
\frac{\partial c_3}{\partial t}
  &= \phantom{D_T\nabla_s^2 c_T +{}} s_3(\mathbf x) - c_3/\tau_3
   \qquad\text{\emph{no diffusion term, and that is the whole point of it}} \label{eq:timp3}
\end{align}
%
with the activation rate --- the ternary mechanism of \S\,``TIMP-2 is the adaptor'' --- carrying one
\emph{occupied} and one \emph{free} MT1-MMP:
%
\begin{equation}
a \;=\; \min\!\Bigl(\,k_{\rm act}\,T^2\,c_P\,\frac{x}{(1+x)^2}\,,\;\frac{c_P}{\Delta t}\Bigr),
\qquad x=\frac{c_T+c_3}{K},
\label{eq:act}
\end{equation}
%
the clamp being the statement that a frame cannot activate more zymogen than it has. MT1-MMP has no
equation of its own because it has no transport and no turnover here: $T(\mathbf x)$ is a per-cell
state, set once and carried, which is exactly the modelling claim that a \emph{tethered} enzyme acts
where its cell put it.

\textbf{And the sheet is what they eat.} The two proteolytic arms --- soluble and tethered --- act on
the areal mass of a face, and a face is removed when its density falls through a critical fraction of
the seeded value:
%
\begin{equation}
\frac{\partial m}{\partial t} = -\,k_{\rm deg}\bigl(c_M + s_{\rm MT1}\,T\bigr)\,m,
\qquad
\rho = \frac{m}{A},\qquad
\text{face dies when }\ \rho/\rho_0 < \rho_{\rm crit}.
\label{eq:degrade}
\end{equation}
%
Eq.~\eqref{eq:degrade} is first-order in $m$, so mass decays exponentially under a steady enzyme and
never goes negative; the discontinuity --- the \emph{hole} --- enters only through $\rho_{\rm crit}$,
which is a material failure threshold and not a chemical one.

\textbf{Two numbers decide whether any of this is resolvable, and one of them currently fails.} A
reaction--diffusion field acquires a length $\sqrt{D/k}$, and that length has to be larger than the
element it is solved on. Measured on the sheet at the published rates, $\sqrt{D/k}=0.05$ against an
element size $h=0.076$: the chemistry's own length is \emph{below} the mesh, which is gate G42's
failure and not a tuning matter. It has a visible consequence in \op{06}: an MT1 source narrower than
about $20^\circ$ of arc tears nothing at all, because a spot smaller than $\sqrt{D/k}$ cannot hold a
concentration against diffusion. \emph{The smallest hole this model can open is set by a diffusion
length, not by a parameter.}

\textbf{The numerics, and why only one of these operators is expensive.} Explicit diffusion needs
$\Delta t\,D/h^2<\tfrac12$ --- $1{,}300$--$13{,}400$ substeps per frame at these $D$, quadrupling at
every $1\!\to\!4$ refinement --- so Eqs.~\eqref{eq:pro}--\eqref{eq:timp} are stepped semi-implicitly,
$(\mathbf I+\Delta t\,D\,\mathbf L)\,c = c_{\rm old}$, by conjugate gradient on a finite-volume face
Laplacian with cotangent weights and face areas. That is unconditionally stable, one solve per frame,
and indifferent to $h$. Eq.~\eqref{eq:timp3} needs no solve at all, which is the computational shadow
of its biology: TIMP-3 is bound to the matrix, so it is a state and not a field.

\textbf{Three honest defects in the discretisation, none of them cosmetic.} (i) In the inhibition
reaction the immobile $c_3$ enters the rate but is not consumed by it --- only $c_M$ and $c_T$ are
decremented --- so TIMP-3 acts as a non-depleting site; that is defensible for a matrix-bound
inhibitor in excess and it is \emph{not} what Eq.~\eqref{eq:timp} says, which is why it is written
here. (ii) The reaction is clamped against $c_M$ alone, so $c_T$ can be driven negative within a frame
and is clipped at zero afterwards: the pair is not conserved at large $k_i\Delta t$. (iii) The
per-frame operator splitting (secrete, diffuse, activate, inhibit, clear) is first-order, so all rates
are per \emph{frame} and the box carries no seconds --- $D$ is converted once from $\mu$m$^2$/s
through the run's declared \code{units:}, and no result may be quoted in seconds without it.

\subsection*{3.3\; Myosin is an ODE, and saying so is the point}

The tissue's contractility is written per \emph{junction}, and it is worth being exact about its
mathematical type because the temptation is to call everything a PDE. Myosin here has \textbf{no
spatial operator}: there is no $\nabla^2$, no flux between neighbouring junctions, and therefore no
length scale of its own. It is a first-order relaxation toward a mechanically-set point,
%
\begin{equation}
\frac{\mathrm d\,\mathrm{myo}_e}{\mathrm dt}
  = \frac{1}{\tau}\Bigl(\mathrm{myo}^{\rm ss}_e - \mathrm{myo}_e\Bigr),
\qquad
\mathrm{myo}^{\rm ss}_e = \alpha\Bigl[\,1 + \sigma\beta\Bigl(\frac{g_e}{\langle g\rangle}-1\Bigr)\Bigr]_+,
\qquad \mathrm{myo}_e\in[0,5],
\label{eq:myosin}
\end{equation}
%
one equation per live junction $e$, where $\alpha$ is the global activity (the in-silico
blebbistatin), $\tau$ the recruitment time, $g_e$ the drive --- edge length, line tension
$\Lambda\,\mathrm{myo}_e\,\ell_e$, or strain rate, chosen by \code{keyed\_on} --- and $\langle
g\rangle$ its \emph{current} mean over live junctions. The sign $\sigma=+1$ is the destabilising
choice: on a tension-keyed drive it is positive feedback, tauter $\to$ more myosin $\to$ tauter, which
is the germband instability and raises the T1 rate; $\sigma=-1$ is tension homeostasis and is kept as
the contrast.

Three details in Eq.~\eqref{eq:myosin} are load-bearing. The reference $\langle g\rangle$ is the
running mean and not a constant, because the tissue triples in radius and a fixed reference would call
every junction ``stretched'' by the end --- myosin would then rise everywhere for a reason that is
growth rather than tension. The activity $\alpha$ multiplies the \emph{set point} rather than the
tension, so inhibition takes effect over $\tau$ the way a drug does instead of instantly. And the
state is keyed to the junction's identity, so it survives a T1, a division and a cell death by
construction rather than by bookkeeping: junctions that no longer exist are simply not written back.

\textbf{So the model reaches the molecule in two different senses, and they are not equally strong.}
The proteases are a genuine reaction--diffusion system on the deformed surface, with named species,
measured dissociation constants and a rate law that predicts a bell rather than reproducing one.
Myosin is a lumped per-junction activity with a mechanosensitive set point --- the right level for a
vertex model, and one equation short of a molecular one, since it has no explicit motor binding, no
phosphorylation state and no spatial transport. Both are written here at the same level of detail so
the difference is visible instead of implied.
